\section{PPO：Proximal Policy Optimization}

\subsection{核心思想：限制策略偏移}

两种策略：

\textbf{方法 1：KL 约束}

$$\max_{\theta'} \tilde{J}(\theta') \quad \text{s.t.} \quad \mathbb{E}_{s \sim \pi_\theta}\left[D_{\mathrm{KL}}\left(\pi_{\theta'}(\cdot \mid s) \| \pi_\theta(\cdot \mid s)\right)\right] \leq \epsilon$$

在 LLM 的偏好优化中非常常见。

\textbf{方法 2：Clipping}

不直接约束策略，而是裁剪 importance weights，移除策略大幅偏离的激励：

$$\tilde{J}(\theta') = \sum_{t,i} \min\left(\frac{\pi_{\theta'}(a_{i,t} \mid s_{i,t})}{\pi_\theta(a_{i,t} \mid s_{i,t})} \hat{A}^{\pi_\theta}, \; \mathrm{clip}\left(\frac{\pi_{\theta'}(a_{i,t} \mid s_{i,t})}{\pi_\theta(a_{i,t} \mid s_{i,t})}, 1-\epsilon, 1+\epsilon\right) \hat{A}^{\pi_\theta}\right)$$

\begin{importantbox}{PPO 的核心目标}
这就是 PPO 的最终 surrogate objective。通过 clipping 和 min 操作，策略不再有激励大幅偏离旧策略，即使做多步更新也能保持稳定。
\end{importantbox}

\subsection{PPO 完整算法}

\begin{enumerate}
    \item 从当前策略 $\pi_\theta$ 采集一批数据
    \item 拟合价值函数 $\hat{V}^\pi_\phi$
    \item 计算优势 $\hat{A}^{\pi_\theta}$
    \item 在这批数据上对 PPO objective 做多步梯度更新（约 3--10 步）
    \item 重复
\end{enumerate}

\subsection{本章小结}

PPO 通过 clipping importance weights 来限制策略偏移，允许在同一批数据上做多步更新，是目前最广泛使用的 on-policy RL 算法之一。

